\documentclass{article}
\usepackage[utf8]{inputenc}
\usepackage{amsmath, amssymb}

\begin{document}
	
	\section*{Problema 1: Pensamiento Matemático - Psicología}
	
	Dadas las proposiciones:
	
	
	\[
	\begin{aligned}
		p &: \text{La memoria a corto plazo tiene capacidad ilimitada.} \\
		q &: \text{La memoria sensorial tiene una duración de varios minutos.} \\
		r &: \text{La memoria a largo plazo es de duración prolongada.} \\
		s &: \text{El individuo utiliza la información para tomar decisiones.}
	\end{aligned}
	\]
	
	
	
	\subsection*{a) Traducción al lenguaje usual}
	
	La proposición compuesta es:
	
	
	\[
	(p \to q) \land (\sim r \to s)
	\]
	
	
	
	Traducción:  
	\textit{“Si la memoria a corto plazo tiene capacidad ilimitada, entonces la memoria sensorial dura varios minutos, y si la memoria a largo plazo no es prolongada, entonces el individuo utiliza la información para tomar decisiones.”}
	
	\subsection*{b) Valor de verdad}
	
	La proposición es:
	
	
	\[
	[(p \land \sim q) \lor r] \to s
	\]
	
	
	
	\textbf{Análisis lógico:}
	
	
	\[
	\begin{aligned}
		p &= \text{Falso} \\
		q &= \text{Falso} \\
		r &= \text{Verdadero} \\
		s &= \text{Verdadero}
	\end{aligned}
	\]
	
	
	
	Evaluación paso a paso:
	
	
	\[
	\sim q = \text{Verdadero}
	\]
	
	
	
	
	\[
	p \land \sim q = \text{Falso} \land \text{Verdadero} = \text{Falso}
	\]
	
	
	
	
	\[
	(p \land \sim q) \lor r = \text{Falso} \lor \text{Verdadero} = \text{Verdadero}
	\]
	
	
	
	
	\[
	[(p \land \sim q) \lor r] \to s = \text{Verdadero} \to \text{Verdadero} = \text{Verdadero}
	\]
	
	
	
	\subsection*{Conclusión}
	El valor de verdad de la proposición compuesta es:
	
	
	\[
	\boxed{\text{Verdadero}}
	\]
	
	
	
\end{document}
